% Short running form; the full title is too wide for the head's right box.
\runninghead{The Dying Son and the Two Peoples}
\properlane{israel-and-the-nations}{The Dying Son and the Two Peoples: Israel's Confession and the Nations' Hope}
\label{sec:israel-and-the-nations}

The Fathers read the healing at Capharnaum not only as one man's faith but as
a figure of peoples. St Augustine reads it in its sequence. The Samaritans
believed Christ's word without a sign; his own country, where he was brought up,
saw signs and scarcely believed; and in the order of these events Augustine
reads the order of salvation, the people of the Jews, among whom Christ was born
and worked his signs, and the nations, who have believed his word having seen
none. Yet the root, he insists, remains in honour: the patriarchs, from whom
Christ came. St Thomas Aquinas reads the ruler himself allegorically, as
Abraham or one of the Fathers of the Old Testament, and the sick son as the
people of Israel, for whom the Fathers prayed. St Cyril of Alexandria reads the
same sequence as the readiness of the nations to obey, and ends with St Paul's
hope that when the fullness of the Gentiles has come in, all Israel will be
saved. Read with that Gospel, the formulary's chants let both peoples be heard.
Israel's own confession in exile opens the Mass, from Azarias's prayer, and
Israel's lament by Babylon's rivers opens the offering; beside them stand the
universal hope of the Gradual, in which the eyes of all look to God for food,
and the Communion's hope in a promised word. The Church prays Israel's
confession as her own, as Azarias prayed for a people's sin he had not
committed; she gives thanks as one of the nations fed; and with Aquinas's
Fathers she prays for the healing of the whole house.

\subsection{Samaria, Galilee and the prophet's country: St Augustine}

Augustine's sixteenth \work{Tractate on John} begins from a puzzle in the text.
After two days among the Samaritans Jesus went into Galilee, ``For Jesus himself
gave testimony that a prophet hath no honour in his own country'' (Jn 4:44).
Why say this of a journey from Samaria, where he was honoured, to Galilee, where
he was brought up? Augustine's answer reads the whole sequence. ``In Samaria He
spent two days, and the Samaritans believed on Him; many were the days He spent
in Galilee, and yet the Galileans did not believe on Him.'' At Cana, even after
the water was made wine, only his disciples believed; now the ruler comes, and
Christ answers him, ``Except ye see signs and wonders, ye believe not''; and
after the healing, ``there did not believe but `himself and his house'.'' ``At
His discourse alone many of the Samaritans believed; at that miracle, in the
place where it was wrought, only that house believed'' (16, 3).

Then Augustine reads the events as signs. The word for a wonder,
\latin{prodigium}, is so called ``as if it were \latin{porrodicium}, \latin{quod
porro dicat}, what betokens something to come'', and these circumstances
portended something. ``Let us, I say, assume the Lord's own country to mean the
people of the Jews,'' among whom he walked and worked, and ``after all that
scarcely a few believed''. Then he turns to his own hearers: ``I am speaking to
God's people; so many of us have believed, what signs have we seen? \dots\ We
have heard the gospel, have given it our consent, have believed on Christ through the
gospel; we have seen no signs, none do we demand'' (16, 3). St Thomas the
Apostle, ``an Israelite, of the Lord's nation'', was rebuked as the ruler was,
and the blessing went to others: ``Blessed are they that have not seen, and yet
have believed. We are the persons here foretold'' (16, 4).

The centurion completes the figure. He ``declared himself unworthy of His
presence''; the ruler ``sought His presence by force''. ``Here is a ceding to
loftiness; there, a conceding to humility.'' ``The centurion, an alien, believed
me able to work by a word, and believed before I did it.'' From this Augustine
draws St Paul's image of the olive: the proud branches broken off, the humble
wild olive grafted in; ``nevertheless let the root remain, while those are cut
off and these received in their place. Where does the root remain? In the
patriarchs. For the people Israel is Christ's own country, since it is of them
that He came according to the flesh; but the root of this tree is Abraham,
Isaac, and Jacob, the holy patriarchs. And where are they? In rest with God, in
great honor'' (16, 5). He ends with the country Christ founded: ``My mother is
Sion, shall a man say'' (Ps 86:5, as Augustine reads it), for from Sion he took flesh, of the Virgin
Mary; ``He had no honor in His country, wherein He was formed; let Him have
honor in the country which He has formed'' (16, 7).

\subsection{The Fathers praying for a sick people: St Thomas Aquinas}

Aquinas, after the literal exposition, gives an allegory that sets the ruler on
the other side. The ruler is Abraham, or one of the Fathers of the Old
Testament, so called because he adheres by faith to the great King, Christ; his
son is the people of the Jews, the seed of Abraham. The son is sick with
depraved pleasures and teachings, and he lies sick in Capharnaum, whose name
Aquinas renders ``abundance'', the abundance that led Israel away from God, as
Moses sang, ``The beloved grew fat, and kicked \dots\ he forsook God who made
him'' (Deut 32:15). The ruler's prayer is the Fathers' intercession: so too the
Fathers of the Old Testament prayed for the people of Israel, as Jeremias is
called ``a lover of his brethren, and of the people of Israel: this is he that
prayeth much for the people'' (2 Mac 15:14). The cure comes by the word at the
seventh hour, and the seventh hour signifies the seven gifts of the Holy Spirit,
through whom sins are forgiven. And Aquinas reads the Galileans' slowness
beside the Samaritans' readiness as the slowness with which, for their
hardness, the Jews came to the faith little by little, \latin{paulatim}, like the
last grapes gleaned after the vintage (Mic 7:1; \work{Super Ioannem} c.~4,
lect.~7).

Augustine and Aquinas agree that the episode figures the relation of Israel and
the nations to Christ's word, and that Israel's patriarchs keep their honour.
They differ on what the ruler figures. For Augustine he belongs to the country
that needed signs; for Aquinas he figures the patriarch who intercedes for that
country. Both stand, each at its own place: Augustine's figure in
the sequence of the chapter, Aquinas's in the person of the father.

\subsection{The readiness of the aliens, and the hope for Israel: St Cyril}

St Cyril of Alexandria reads the same sequence forward into the next chapter,
where the healing at the pool in Jerusalem brings persecution. The evangelist's
aim was ``to shew how superior in obedience were the aliens to the Jews'', and
from the contrast of the chapters, ``thus and in no other way could we learn,
that by the just judgment of God \dots\ Israel with reason falleth from the hope, and the fulness of the Gentiles
is brought in in his place''. ``He had by one miracle saved the city of the
Samaritans, by one likewise the nobleman''; ``the one believed with his whole
house, and confessed that Jesus is God''. But Cyril does not end there. Those
set in the first rank ``because of the election of the fathers, will come last
and after the calling of the Gentiles. For when the fulness of the Gentiles is
come in, then shall all Israel be saved'' (\work{Commentary on John} II.5;
Rom 11:25--26). Augustine's tractate does not say this; it is Cyril's hope.

\subsection{Israel's voice in the Church's mouth: the Introit and the Offertory}

The two chants of exile are, in their first sense, Israel's. The Introit is
Azarias's confession for a people taken captive for its sins; the Offertory is
the lament of that people by Babylon's rivers. St Jerome hears in Azarias's
\latin{peccavimus} the innocent speaking \latin{ex persona populi}, in the
person of the people, and he gives the Church leave to pray the same prayer
when, for the people's sins, she suffers want (\work{Commentarii in Danielem},
on 3:29, 37). The Church therefore prays Israel's confession as her own, and
prays it as Azarias did, for a people. St John Chrysostom, expounding the
Offertory's psalm of the historical captives, describes what the exile taught
them: they had scorned the city while they held it and mocked the prophets who
wept for it, and in Babylon they wept unbidden, kept the law among their
captors, and refused their songs to profane ears (\work{Expositio in Psalmos},
on Ps~136).

The medieval commentators on these chants hear the same voices. Rupert of
Deutz calls the Introit the voice of the prophets and reads the Offertory's
harps, hung on the willows, as preachers fallen silent before barren hearers
(\work{De divinis officiis} XII.20). Honorius Augustodunensis expounds the
chants \latin{sub lege}, as the song of the exiles who returned from Babylon
(\work{Gemma animae} IV.88--89). Dom Lucien Fromage, in the continuation of
\work{The Liturgical Year}, gives the chants to two choirs: ``The Gentile-Choir
takes the Gradual and Communion-Anthem; the Choir of the Jews, the Introit and
Offertory''. ``In the Introit, the Jewish people sang its repentance and humble
confidence''; in the Offertory Israel, grown repentant, ``joins our Mother the
Church'' in singing the psalm of exile (vol.~XI, pp.~439, 447, 451). The Introit's verse gives one further voice. Cassiodorus hears the whole of
Ps~118 spoken by the single chorus of the saints, from the world's beginning to
its end, set forth as one speaker for the unity of the Church
(\work{Expositio Psalmorum} 118, preface); the blessing of those who walk in
the law of the Lord is then spoken by the saints of both peoples together.

\subsection{All eyes and every living thing: the nations' chants and the orations}

The Gradual is universal in its words: \latin{Oculi omnium}, the eyes of all,
and \latin{omne animal}, every living creature, filled with blessing. Fromage
gives it to the Gentiles. The Alleluia's verse continues, in the psalm, to the
nations: ``I will praise thee, O Lord, among the people: and I will sing unto
thee among the nations'' (Ps 107:4). Cassiodorus reads that confession among the
peoples and that song among the nations as the Church gathered first from the
Jews and then from the Gentiles (\work{Expositio Psalmorum} 107.4). The
Communion's hope rests on a promise, and St Robert Bellarmine names it: ``thy
promise `in which thou hast given me hope,' when you said to Abraham, and
through him to all his children, `Walk before me, and be perfect;' I will be
`thy reward exceedingly great'.'' St Augustine reads the Communion's comfort as
the hope given ``to Christ's Body, that is, to the Church''. The Church of the
nations hopes in the word first given to Abraham.

The Epistle is the nations' new life, spoken to Gentile Ephesians who were
``heretofore darkness, but now light in the Lord'' (Eph 5:8): walk as wise,
redeem the time, be filled with the Spirit, give thanks. The Collect asks pardon and peace for \latin{fidelibus tuis}, the
faithful of every people grafted in. The Secret asks heavenly medicine for the
heart's vices, the cure Aquinas's sick son needed for the depraved pleasures that
abundance bred. And the Postcommunion asks obedience to the commandments, the
law given to Israel and confessed broken in the Introit, now asked of all who
receive the sacred gifts.

\subsection{Whose story? The difficulties of the reading}

The reading can be heard as setting Israel against the Church. Its witnesses
say both more and less than that. Augustine insists that the root, the
patriarchs, remains in honour, and that Christ took flesh from Sion, which he
calls his mother; Aquinas's ruler is Abraham interceding for his people; Cyril
ends in the salvation of all Israel; Jerome makes Israel's confession the
prayer of the innocent and the Church's own. The figure the Fathers read is a
figure of the history of salvation, told in their own terms, and the honour
Augustine gives the root and the hope with which Cyril ends belong to it as
much as the broken branches. Augustine's own conclusion is humility for
the grafted branch, since it stands by faith and not by birth; Aquinas's, the
intercession of the Fathers for their people; Cyril's, hope.

The identifications do not agree.
Augustine's ruler belongs to the people that needed signs; Aquinas's figures
the patriarch who intercedes for it; Cyril sets him beside the Samaritans among
the obedient aliens. What the three share is the controlling claim: that the
healing by the word figures the relation of Israel and the nations to that
word.

The chants raise a last difficulty. Fromage's two choirs give the Introit and
Offertory to Israel; St Augustine, expounding the Offertory's psalm, hears in it
the Church's own captivity, the voice the second reading hears. Both are
traditional answers, and they differ on whose exile the chant first voices. Israel's
voice in the chants is Fromage's answer and the commentators'; the figure in the
Gospel is Augustine's own; and Jerome's \latin{ex persona populi} joins the two, since the Church, praying Israel's confession, prays as Azarias did,
for and with a people.

\subsection{The four senses of this reading}

\begin{fourSenses}
\item[Literal.] A royal official asks for his son's life at Cana after the
Samaritans have believed Jesus's word without a sign and the Galileans have
received him for what they saw at Jerusalem (Jn 4:39--54); in the Introit
Azarias confesses for exiled Israel (Dan 3:26--45), and in the Offertory the
captives lament by Babylon's rivers (Ps 136:1).

\item[Allegorical.] Christ's own country figures the people of the Jews, who saw
signs and scarcely believed, and the Samaritans the nations, who believed the
word without signs (Augustine); the ruler figures the Fathers of the Old
Testament interceding, and the son their people, healed by the word at the hour
of the sevenfold Spirit (Aquinas); the nations' psalm among the peoples is the
Church of the Jews first and then of the Gentiles (Cassiodorus).

\item[Moral.] The Church of the nations is the grafted branch, and must be humble
before the root, the patriarchs in great honour (Augustine); she prays Israel's
confession as her own, as the innocent prayed it for the people (Jerome), and
intercedes as the Fathers interceded (Aquinas).

\item[Anagogical.] ``When the fulness of the Gentiles is come in, then shall all
Israel be saved'' (Cyril, after Rom 11:25--26): the healing of the whole house
at Capharnaum looks to the healing of both peoples in the one house of God.
\end{fourSenses}
